% Options for packages loaded elsewhere
\PassOptionsToPackage{unicode}{hyperref}
\PassOptionsToPackage{hyphens}{url}
\documentclass[
]{article}
\usepackage{xcolor}
\usepackage{amsmath,amssymb}
\setcounter{secnumdepth}{-\maxdimen} % remove section numbering
\usepackage{iftex}
\ifPDFTeX
  \usepackage[T1]{fontenc}
  \usepackage[utf8]{inputenc}
  \usepackage{textcomp} % provide euro and other symbols
\else % if luatex or xetex
  \usepackage{unicode-math} % this also loads fontspec
  \defaultfontfeatures{Scale=MatchLowercase}
  \defaultfontfeatures[\rmfamily]{Ligatures=TeX,Scale=1}
\fi
\usepackage{lmodern}
\ifPDFTeX\else
  % xetex/luatex font selection
\fi
% Use upquote if available, for straight quotes in verbatim environments
\IfFileExists{upquote.sty}{\usepackage{upquote}}{}
\IfFileExists{microtype.sty}{% use microtype if available
  \usepackage[]{microtype}
  \UseMicrotypeSet[protrusion]{basicmath} % disable protrusion for tt fonts
}{}
\makeatletter
\@ifundefined{KOMAClassName}{% if non-KOMA class
  \IfFileExists{parskip.sty}{%
    \usepackage{parskip}
  }{% else
    \setlength{\parindent}{0pt}
    \setlength{\parskip}{6pt plus 2pt minus 1pt}}
}{% if KOMA class
  \KOMAoptions{parskip=half}}
\makeatother
\usepackage{color}
\usepackage{fancyvrb}
\newcommand{\VerbBar}{|}
\newcommand{\VERB}{\Verb[commandchars=\\\{\}]}
\DefineVerbatimEnvironment{Highlighting}{Verbatim}{commandchars=\\\{\}}
% Add ',fontsize=\small' for more characters per line
\newenvironment{Shaded}{}{}
\newcommand{\AlertTok}[1]{\textcolor[rgb]{1.00,0.00,0.00}{\textbf{#1}}}
\newcommand{\AnnotationTok}[1]{\textcolor[rgb]{0.38,0.63,0.69}{\textbf{\textit{#1}}}}
\newcommand{\AttributeTok}[1]{\textcolor[rgb]{0.49,0.56,0.16}{#1}}
\newcommand{\BaseNTok}[1]{\textcolor[rgb]{0.25,0.63,0.44}{#1}}
\newcommand{\BuiltInTok}[1]{\textcolor[rgb]{0.00,0.50,0.00}{#1}}
\newcommand{\CharTok}[1]{\textcolor[rgb]{0.25,0.44,0.63}{#1}}
\newcommand{\CommentTok}[1]{\textcolor[rgb]{0.38,0.63,0.69}{\textit{#1}}}
\newcommand{\CommentVarTok}[1]{\textcolor[rgb]{0.38,0.63,0.69}{\textbf{\textit{#1}}}}
\newcommand{\ConstantTok}[1]{\textcolor[rgb]{0.53,0.00,0.00}{#1}}
\newcommand{\ControlFlowTok}[1]{\textcolor[rgb]{0.00,0.44,0.13}{\textbf{#1}}}
\newcommand{\DataTypeTok}[1]{\textcolor[rgb]{0.56,0.13,0.00}{#1}}
\newcommand{\DecValTok}[1]{\textcolor[rgb]{0.25,0.63,0.44}{#1}}
\newcommand{\DocumentationTok}[1]{\textcolor[rgb]{0.73,0.13,0.13}{\textit{#1}}}
\newcommand{\ErrorTok}[1]{\textcolor[rgb]{1.00,0.00,0.00}{\textbf{#1}}}
\newcommand{\ExtensionTok}[1]{#1}
\newcommand{\FloatTok}[1]{\textcolor[rgb]{0.25,0.63,0.44}{#1}}
\newcommand{\FunctionTok}[1]{\textcolor[rgb]{0.02,0.16,0.49}{#1}}
\newcommand{\ImportTok}[1]{\textcolor[rgb]{0.00,0.50,0.00}{\textbf{#1}}}
\newcommand{\InformationTok}[1]{\textcolor[rgb]{0.38,0.63,0.69}{\textbf{\textit{#1}}}}
\newcommand{\KeywordTok}[1]{\textcolor[rgb]{0.00,0.44,0.13}{\textbf{#1}}}
\newcommand{\NormalTok}[1]{#1}
\newcommand{\OperatorTok}[1]{\textcolor[rgb]{0.40,0.40,0.40}{#1}}
\newcommand{\OtherTok}[1]{\textcolor[rgb]{0.00,0.44,0.13}{#1}}
\newcommand{\PreprocessorTok}[1]{\textcolor[rgb]{0.74,0.48,0.00}{#1}}
\newcommand{\RegionMarkerTok}[1]{#1}
\newcommand{\SpecialCharTok}[1]{\textcolor[rgb]{0.25,0.44,0.63}{#1}}
\newcommand{\SpecialStringTok}[1]{\textcolor[rgb]{0.73,0.40,0.53}{#1}}
\newcommand{\StringTok}[1]{\textcolor[rgb]{0.25,0.44,0.63}{#1}}
\newcommand{\VariableTok}[1]{\textcolor[rgb]{0.10,0.09,0.49}{#1}}
\newcommand{\VerbatimStringTok}[1]{\textcolor[rgb]{0.25,0.44,0.63}{#1}}
\newcommand{\WarningTok}[1]{\textcolor[rgb]{0.38,0.63,0.69}{\textbf{\textit{#1}}}}
\usepackage{longtable,booktabs,array}
\usepackage{caption}
\captionsetup[table]{skip=6pt}
\newcounter{none} % for unnumbered tables
\usepackage{calc} % for calculating minipage widths
% Correct order of tables after \paragraph or \subparagraph
\usepackage{etoolbox}
\makeatletter
\patchcmd\longtable{\par}{\if@noskipsec\mbox{}\fi\par}{}{}
\makeatother
\setlength{\emergencystretch}{3em} % prevent overfull lines
\providecommand{\tightlist}{%
  \setlength{\itemsep}{0pt}\setlength{\parskip}{0pt}}
\usepackage{bookmark}
\IfFileExists{xurl.sty}{\usepackage{xurl}}{} % add URL line breaks if available
\urlstyle{same}
% fallback for those not using the hyperref driver hyperxmp:
\makeatletter
\@ifundefined{xmpquote}{\newcommand{\xmpquote}[1]{#1}}{}
\makeatother
\hypersetup{
  hidelinks,
  pdfcreator={LaTeX via pandoc}}

\author{}
\date{}

\begin{document}

\section{Tutoriel Pratique Exhaustif COVOF 2D (C++17 /
OpenMPI)}\label{tutoriel-pratique-exhaustif-covof-2d-c17-openmpi}

Ce document constitue le \textbf{guide d'apprentissage complet pas à
pas} pour maîtriser l'utilisation, la configuration, l'exécution et
l'analyse de simulations numériques avec le solveur CFD \textbf{COVOF
2D}.

\begin{center}\rule{0.5\linewidth}{0.5pt}\end{center}

\subsection{Table des Matières}\label{table-des-matiuxe8res}

\begin{enumerate}
\def\labelenumi{\arabic{enumi}.}
\tightlist
\item
  \hyperref[1-introduction-et-notions-fondamentales]{Introduction et
  Notions Fondamentales}
\item
  \hyperref[2-installation-et-environnement-de-compilation]{Installation
  et Environnement de Compilation}
\item
  \hyperref[3-structure-et-syntaxe-dun-cas-test-casesjson]{Structure et
  Syntaxe d'un Cas Test (\texttt{cases/*.json})}
\item
  \hyperref[4-tutoriel-1--uxe9coulement-de-poiseuille-et-transitoire-de-szymanski]{Tutoriel
  1 : Écoulement de Poiseuille et Transitoire de Szymanski}
\item
  \hyperref[5-tutoriel-2--cavituxe9-entrauxeenuxe9e-uxe0-reynolds-re--1000]{Tutoriel
  2 : Cavité Entraînée à Reynolds Re = 1000}
\item
  \hyperref[6-tutoriel-3--uxe9coulement-sur-marche-descendante-armaly-re--100]{Tutoriel
  3 : Écoulement sur Marche Descendante (Armaly Re = 100)}
\item
  \hyperref[7-tutoriel-4--convection-naturelle-thermique-de-vahl-davis-ra--10ux2074]{Tutoriel
  4 : Convection Naturelle Thermique de Vahl Davis (Ra = 10⁴)}
\item
  \hyperref[8-tutoriel-5--uxe9coulement-diphasique-vof--rupture-de-barrage]{Tutoriel
  5 : Écoulement Diphasique VOF --- Rupture de Barrage}
\item
  \hyperref[9-tutoriel-6--tension-de-surface-et-saut-de-pression-de-laplace]{Tutoriel
  6 : Tension de Surface et Saut de Pression de Laplace}
\item
  \hyperref[10-tutoriel-7--uxe9coulement-autour-dun-cylindre-par-puxe9nalisation-de-brinkman]{Tutoriel
  7 : Écoulement Autour d'un Cylindre par Pénalisation de Brinkman}
\item
  \hyperref[11-post-traitement-avancuxe9-avec-paraview]{Post-Traitement
  Avancé avec ParaView}
\item
  \hyperref[12-automatisation-python-et-uxe9tudes-de-convergence]{Automatisation
  Python et Études de Convergence}
\item
  \hyperref[13-cruxe9er-un-nouveau-cas-test-uxe0-partir-de-zuxe9ro]{Créer
  un Nouveau Cas Test à Partir de Zéro}
\end{enumerate}

\begin{center}\rule{0.5\linewidth}{0.5pt}\end{center}

\subsection{1. Introduction et Notions
Fondamentales}\label{introduction-et-notions-fondamentales}

\subsubsection{1.1. Pourquoi la grille décalée (Staggered MAC Grid)
?}\label{pourquoi-la-grille-duxe9caluxe9e-staggered-mac-grid}

Dans les méthodes traditionnelles à variables colocalisées (où vitesses
et pressions sont aux mêmes nœuds), la discrétisation par différences
finies centrées d'ordre 2 conduit au phénomène bien connu de
\textbf{découplage pair-impair} (modes en damier ou \emph{checkerboard
modes}). Pour éliminer ce problème sans introduire de dissipation
artificielle (comme la projection de Rhie-Chow), COVOF 2D utilise la
grille de Harlow \& Welch (1965) :

\begin{verbatim}
                  v(i, j+1/2)
            +----------▲----------+
            |          |          |
            |          |          |
u(i-1/2, j) ►     p(i, j)         ► u(i+1/2, j)
            |     T(i, j)         |
            |     C(i, j)         |
            |          |          |
            +----------▲----------+
                  v(i, j-1/2)
\end{verbatim}

\begin{itemize}
\tightlist
\item
  \textbf{Vitesse horizontale \(u\)} : située aux faces verticales
  \(x_{i-1/2} = (i-1)\Delta x\).
\item
  \textbf{Vitesse verticale \(v\)} : située aux faces horizontales
  \(y_{j-1/2} = (j-1)\Delta y\).
\item
  \textbf{Scalaires \(p, T, C\)} : situés au barycentre des mailles
  \(x_i = (i-1/2)\Delta x, y_j = (j-1/2)\Delta y\).
\end{itemize}

\subsubsection{1.2. Algorithme de résolution (Lagrangien
Augmenté)}\label{algorithme-de-ruxe9solution-lagrangien-augmentuxe9}

L'équation de Navier-Stokes incompressible s'écrit :
\[\rho \left( \frac{\partial \mathbf{u}}{\partial t} + \mathbf{u} \cdot \nabla \mathbf{u} \right) - \nabla \cdot \left( \mu (\nabla \mathbf{u} + \nabla^T \mathbf{u}) \right) + \nabla p = \rho \mathbf{g} + \mathbf{f}_{\sigma}\]
avec \(\nabla \cdot \mathbf{u} = 0\).

Le Lagrangien Augmenté consiste à ajouter un terme de pénalité
\(-r_{\text{al}} \nabla (\nabla \cdot \mathbf{u})\) à l'équation de
quantité de mouvement, couplé à une mise à jour d'Uzawa :
\[p^{m+1} = p^m - r_{\text{al}} \nabla \cdot \mathbf{u}^{m+1}\] Cette
technique permet de forcer la divergence à zéro dès
\(r_{\text{al}} \sim 100\), tout en garantissant un excellent
conditionnement du système linéaire.

\begin{center}\rule{0.5\linewidth}{0.5pt}\end{center}

\subsection{2. Installation et Environnement de
Compilation}\label{installation-et-environnement-de-compilation}

\subsubsection{2.1. Vérification de
l'environnement}\label{vuxe9rification-de-lenvironnement}

Vérifiez que votre système dispose d'un compilateur supportant C++17 et
d'une distribution OpenMPI :

\begin{Shaded}
\begin{Highlighting}[]
\ExtensionTok{mpic++} \AttributeTok{{-}{-}version}
\FunctionTok{cmake} \AttributeTok{{-}{-}version}
\end{Highlighting}
\end{Shaded}

\subsubsection{2.2. Compilation}\label{compilation}

Depuis la racine du dépôt \texttt{covof2d} :

\begin{Shaded}
\begin{Highlighting}[]
\CommentTok{\# Configuration hors{-}source dans build/}
\FunctionTok{cmake} \AttributeTok{{-}B}\NormalTok{ build }\AttributeTok{{-}DCMAKE\_BUILD\_TYPE}\OperatorTok{=}\NormalTok{Release}

\CommentTok{\# Compilation parallèle sur 4 cœurs}
\FunctionTok{cmake} \AttributeTok{{-}{-}build}\NormalTok{ build }\AttributeTok{{-}j4}
\end{Highlighting}
\end{Shaded}

L'exécutable généré est produit dans \texttt{bin/covof2d}.

\begin{center}\rule{0.5\linewidth}{0.5pt}\end{center}

\subsection{\texorpdfstring{3. Structure et Syntaxe d'un Cas Test
(\texttt{cases/*.json})}{3. Structure et Syntaxe d'un Cas Test (cases/*.json)}}\label{structure-et-syntaxe-dun-cas-test-cases.json}

Chaque simulation est entièrement définie par un fichier JSON structuré
en 7 blocs indépendants :

\begin{Shaded}
\begin{Highlighting}[]
\FunctionTok{\{}
  \DataTypeTok{"mode"}\FunctionTok{:} \StringTok{"NAVIER\_STOKES"}\FunctionTok{,}
  \DataTypeTok{"mesh"}\FunctionTok{:} \FunctionTok{\{}
    \DataTypeTok{"nx"}\FunctionTok{:} \DecValTok{64}\FunctionTok{,}
    \DataTypeTok{"ny"}\FunctionTok{:} \DecValTok{64}\FunctionTok{,}
    \DataTypeTok{"lx"}\FunctionTok{:} \FloatTok{1.0}\FunctionTok{,}
    \DataTypeTok{"ly"}\FunctionTok{:} \FloatTok{1.0}
  \FunctionTok{\},}
  \DataTypeTok{"physics"}\FunctionTok{:} \FunctionTok{\{}
    \DataTypeTok{"enable\_two\_phase"}\FunctionTok{:} \KeywordTok{false}\FunctionTok{,}
    \DataTypeTok{"enable\_convection"}\FunctionTok{:} \KeywordTok{true}\FunctionTok{,}
    \DataTypeTok{"enable\_energy"}\FunctionTok{:} \KeywordTok{false}\FunctionTok{,}
    \DataTypeTok{"gx"}\FunctionTok{:} \FloatTok{0.0}\FunctionTok{,}
    \DataTypeTok{"gy"}\FunctionTok{:} \FloatTok{0.0}\FunctionTok{,}
    \DataTypeTok{"fluid1"}\FunctionTok{:} \FunctionTok{\{} \DataTypeTok{"rho"}\FunctionTok{:} \FloatTok{1.0}\FunctionTok{,} \DataTypeTok{"mu"}\FunctionTok{:} \FloatTok{0.001} \FunctionTok{\}}
  \FunctionTok{\},}
  \DataTypeTok{"time"}\FunctionTok{:} \FunctionTok{\{}
    \DataTypeTok{"dt"}\FunctionTok{:} \FloatTok{0.005}\FunctionTok{,}
    \DataTypeTok{"t\_final"}\FunctionTok{:} \FloatTok{5.0}\FunctionTok{,}
    \DataTypeTok{"cfl\_max"}\FunctionTok{:} \FloatTok{0.5}\FunctionTok{,}
    \DataTypeTok{"adaptive\_dt"}\FunctionTok{:} \KeywordTok{false}
  \FunctionTok{\},}
  \DataTypeTok{"boundary\_conditions"}\FunctionTok{:} \FunctionTok{\{}
    \DataTypeTok{"left"}\FunctionTok{:}   \FunctionTok{\{} \DataTypeTok{"type"}\FunctionTok{:} \StringTok{"DIRICHLET"}\FunctionTok{,} \DataTypeTok{"u"}\FunctionTok{:} \FloatTok{0.0}\FunctionTok{,} \DataTypeTok{"v"}\FunctionTok{:} \FloatTok{0.0} \FunctionTok{\},}
    \DataTypeTok{"right"}\FunctionTok{:}  \FunctionTok{\{} \DataTypeTok{"type"}\FunctionTok{:} \StringTok{"DIRICHLET"}\FunctionTok{,} \DataTypeTok{"u"}\FunctionTok{:} \FloatTok{0.0}\FunctionTok{,} \DataTypeTok{"v"}\FunctionTok{:} \FloatTok{0.0} \FunctionTok{\},}
    \DataTypeTok{"bottom"}\FunctionTok{:} \FunctionTok{\{} \DataTypeTok{"type"}\FunctionTok{:} \StringTok{"DIRICHLET"}\FunctionTok{,} \DataTypeTok{"u"}\FunctionTok{:} \FloatTok{0.0}\FunctionTok{,} \DataTypeTok{"v"}\FunctionTok{:} \FloatTok{0.0} \FunctionTok{\},}
    \DataTypeTok{"top"}\FunctionTok{:}    \FunctionTok{\{} \DataTypeTok{"type"}\FunctionTok{:} \StringTok{"DIRICHLET"}\FunctionTok{,} \DataTypeTok{"u"}\FunctionTok{:} \FloatTok{1.0}\FunctionTok{,} \DataTypeTok{"v"}\FunctionTok{:} \FloatTok{0.0} \FunctionTok{\}}
  \FunctionTok{\},}
  \DataTypeTok{"solver"}\FunctionTok{:} \FunctionTok{\{}
    \DataTypeTok{"type"}\FunctionTok{:} \StringTok{"AUGMENTED\_LAGRANGIAN"}\FunctionTok{,}
    \DataTypeTok{"tol"}\FunctionTok{:} \FloatTok{1.0e{-}7}\FunctionTok{,}
    \DataTypeTok{"max\_iter"}\FunctionTok{:} \DecValTok{500}\FunctionTok{,}
    \DataTypeTok{"r\_al"}\FunctionTok{:} \FloatTok{100.0}\FunctionTok{,}
    \DataTypeTok{"uzawa\_iter"}\FunctionTok{:} \DecValTok{1}\FunctionTok{,}
    \DataTypeTok{"linear\_solver"}\FunctionTok{:} \StringTok{"BICGSTAB2"}\FunctionTok{,}
    \DataTypeTok{"use\_precond"}\FunctionTok{:} \KeywordTok{true}\FunctionTok{,}
    \DataTypeTok{"precond\_type"}\FunctionTok{:} \StringTok{"JACOBI"}
  \FunctionTok{\},}
  \DataTypeTok{"output"}\FunctionTok{:} \FunctionTok{\{}
    \DataTypeTok{"output\_dir"}\FunctionTok{:} \StringTok{"output\_mon\_cas"}\FunctionTok{,}
    \DataTypeTok{"case\_name"}\FunctionTok{:} \StringTok{"mon\_cas"}\FunctionTok{,}
    \DataTypeTok{"vtk\_interval"}\FunctionTok{:} \DecValTok{20}\FunctionTok{,}
    \DataTypeTok{"log\_interval"}\FunctionTok{:} \DecValTok{10}\FunctionTok{,}
    \DataTypeTok{"write\_csv\_history"}\FunctionTok{:} \KeywordTok{true}
  \FunctionTok{\}}
\FunctionTok{\}}
\end{Highlighting}
\end{Shaded}

\subsubsection{Table des Options
Disponibles}\label{table-des-options-disponibles}

{\def\LTcaptype{none} % do not increment counter
\begin{longtable}[]{@{}
  >{\raggedright\arraybackslash}p{(\linewidth - 6\tabcolsep) * \real{0.2500}}
  >{\raggedright\arraybackslash}p{(\linewidth - 6\tabcolsep) * \real{0.2500}}
  >{\raggedright\arraybackslash}p{(\linewidth - 6\tabcolsep) * \real{0.2500}}
  >{\raggedright\arraybackslash}p{(\linewidth - 6\tabcolsep) * \real{0.2500}}@{}}
\toprule\noalign{}
\begin{minipage}[b]{\linewidth}\raggedright
Bloc
\end{minipage} & \begin{minipage}[b]{\linewidth}\raggedright
Paramètre
\end{minipage} & \begin{minipage}[b]{\linewidth}\raggedright
Types / Valeurs Possibles
\end{minipage} & \begin{minipage}[b]{\linewidth}\raggedright
Description
\end{minipage} \\
\midrule\noalign{}
\endhead
\bottomrule\noalign{}
\endlastfoot
\texttt{mode} & \texttt{mode} & \texttt{"NAVIER\_STOKES"},
\texttt{"HEAT"} & Physique principale résolue \\
\texttt{mesh} & \texttt{nx}, \texttt{ny} & Entiers positifs (\(\ge 16\))
& Résolution spatiale en cellules \\
\texttt{mesh} & \texttt{lx}, \texttt{ly} & Réels en mètres & Dimensions
géométriques du domaine \\
\texttt{physics} & \texttt{enable\_convection} & \texttt{true},
\texttt{false} & Active/désactive les termes convectifs
\((\mathbf{u}\cdot\nabla)\mathbf{u}\) \\
\texttt{physics} & \texttt{enable\_two\_phase} & \texttt{true},
\texttt{false} & Active le transport VOF diphasique \\
\texttt{physics} & \texttt{sigma} & Réel en \(\text{N/m}\) & Tension
interfaciale pour le modèle CSF \\
\texttt{boundary} & \texttt{type} & \texttt{"DIRICHLET"},
\texttt{"NEUMANN"}, \texttt{"PERIODIC"}, \texttt{"FREESLIP"} & Type de
condition aux limites \\
\texttt{solver} & \texttt{linear\_solver} & \texttt{"BICGSTAB2"},
\texttt{"BICGSTAB"}, \texttt{"CG"}, \texttt{"GMRES"} & Algorithme de
Krylov \\
\texttt{solver} & \texttt{r\_al} & Réel (\(\ge 0\)) & Paramètre de
pénalité de divergence \\
\end{longtable}
}

\begin{center}\rule{0.5\linewidth}{0.5pt}\end{center}

\subsection{4. Tutoriel 1 : Écoulement de Poiseuille et Transitoire de
Szymanski}\label{tutoriel-1-uxe9coulement-de-poiseuille-et-transitoire-de-szymanski}

Ce cas teste l'établissement transitoire d'un profil parabolique de
Poiseuille sous l'action d'un gradient axial de pression constant
\(G = 0.0555\text{ Pa/m}\) entre deux plaques distantes de
\(2e = 0.16\text{ m}\).

\subsubsection{\texorpdfstring{4.1. Configuration
(\texttt{cases/poiseuille\_unsteady.json})}{4.1. Configuration (cases/poiseuille\_unsteady.json)}}\label{configuration-casespoiseuille_unsteady.json}

\begin{itemize}
\tightlist
\item
  \textbf{Propriétés} : \(\rho = 1.0\text{ kg/m}^3\),
  \(\mu = 0.01\text{ Pa}\cdot\text{s}\), \(g_x = 0.0555\text{ m/s}^2\).
\item
  \textbf{Conditions limites} : Périodique en \(x\), adhérence
  (\(u=0, v=0\)) en haut et en bas.
\item
  \textbf{Pénalité \(r_{\text{al}}\)} : fixé à \(0.0\) car le domaine
  est périodique et la pression stationnaire est nulle (\(p \equiv 0\)).
\end{itemize}

\subsubsection{4.2. Lancement}\label{lancement}

\begin{Shaded}
\begin{Highlighting}[]
\ExtensionTok{mpirun} \AttributeTok{{-}np}\NormalTok{ 2 bin/covof2d cases/poiseuille\_unsteady.json}
\end{Highlighting}
\end{Shaded}

\subsubsection{4.3. Solution Analytique (Série de
Szymanski)}\label{solution-analytique-suxe9rie-de-szymanski}

La solution transitoire exacte s'exprime par la décomposition de Fourier
:
\[u(z, t) = \frac{G e^2}{2\mu} \left( 1 - \frac{z^2}{e^2} \right) + \sum_{n=0}^{\infty} a_n \exp\left( - \frac{(2n+1)^2 \pi^2 \mu}{4 \rho e^2} t \right) \cos\left( \frac{(2n+1) \pi z}{2e} \right)\]
avec \(a_n = -\frac{16 G e^2 (-1)^n}{\mu \pi^3 (2n+1)^3}\) et
\(z = y - e \in [-e, e]\).

À \(t = 1.0\text{ s}\), la solution COVOF 2D converge à
\(u_{\max} = 0.017343\text{ m/s}\) avec une erreur
\(L_1 = 1.49 \times 10^{-5}\text{ m/s}\), surpassant la cible de
référence Référence Académique (\(3.13 \times 10^{-4}\text{ m/s}\)).

\begin{center}\rule{0.5\linewidth}{0.5pt}\end{center}

\subsection{5. Tutoriel 2 : Cavité Entraînée à Reynolds Re =
1000}\label{tutoriel-2-cavituxe9-entrauxeenuxe9e-uxe0-reynolds-re-1000}

Le benchmark de Ghia et al.~(1982) est le problème classique pour tester
la précision des termes convectifs et visqueux.

\begin{verbatim}
       u_top = 1.0 m/s, v_top = 0.0 m/s
    +--------------------------------+
    |                                |
    |           VORTEX               |
    |          PRIMAIRE              |
    |                                |
    |  Vortex                  Vortex|
    |  coin bas-gauche   coin bas-dr.|
    +--------------------------------+
             Parois fixes (u=0, v=0)
\end{verbatim}

\subsubsection{\texorpdfstring{5.1. Configuration
(\texttt{cases/cavity\_re1000.json})}{5.1. Configuration (cases/cavity\_re1000.json)}}\label{configuration-casescavity_re1000.json}

\begin{itemize}
\tightlist
\item
  Domaine carré \(1.0\text{ m} \times 1.0\text{ m}\), maillage
  \(128 \times 128\).
\item
  Fluide : \(\rho = 1.0\text{ kg/m}^3\),
  \(\mu = 0.001\text{ Pa}\cdot\text{s} \implies Re = \frac{\rho U L}{\mu} = 1000\).
\item
  Paroi supérieure en mouvement à \(u = 1.0\text{ m/s}\).
\end{itemize}

\subsubsection{5.2. Lancement}\label{lancement-1}

\begin{Shaded}
\begin{Highlighting}[]
\ExtensionTok{mpirun} \AttributeTok{{-}np}\NormalTok{ 4 bin/covof2d cases/cavity\_re1000.json}
\end{Highlighting}
\end{Shaded}

\subsubsection{5.3. Résultats attendus}\label{ruxe9sultats-attendus}

Le vortex primaire central se stabilise à
\((x, y) \approx (0.53, 0.56)\) et deux recirculations secondaires
apparaissent dans les coins inférieurs.

\begin{center}\rule{0.5\linewidth}{0.5pt}\end{center}

\subsection{6. Tutoriel 3 : Écoulement sur Marche Descendante (Armaly Re
=
100)}\label{tutoriel-3-uxe9coulement-sur-marche-descendante-armaly-re-100}

L'expérience d'Armaly et al.~(1983) étudie le décollement de couche
limite et la longueur de recirculation \(x_r / H\) en aval d'une marche
de hauteur \(H\).

\subsubsection{\texorpdfstring{6.1. Configuration
(\texttt{cases/backward\_facing\_step.json})}{6.1. Configuration (cases/backward\_facing\_step.json)}}\label{configuration-casesbackward_facing_step.json}

\begin{itemize}
\item
  Longueur \(L_x = 40.0\text{ m}\), Hauteur \(L_y = 2.0\text{ m}\).
\item
  Hauteur de la marche : \(H = 1.0\text{ m}\) sur
  \(x \in [0, 4.0\text{ m}]\).
\item
  Profil d'entrée parabolique sur \([H, L_y]\) avec
  \(U_{\text{inlet}} = 1.0\text{ m/s}\).
\item
  L'obstacle de la marche est modélisé par pénalisation de Brinkman :

\begin{Shaded}
\begin{Highlighting}[]
\StringTok{"obstacle"}\ErrorTok{:} \FunctionTok{\{}
  \DataTypeTok{"enabled"}\FunctionTok{:} \KeywordTok{true}\FunctionTok{,}
  \DataTypeTok{"type"}\FunctionTok{:} \StringTok{"BOX"}\FunctionTok{,}
  \DataTypeTok{"x0"}\FunctionTok{:} \FloatTok{0.0}\FunctionTok{,} \DataTypeTok{"x1"}\FunctionTok{:} \FloatTok{4.0}\FunctionTok{,}
  \DataTypeTok{"y0"}\FunctionTok{:} \FloatTok{0.0}\FunctionTok{,} \DataTypeTok{"y1"}\FunctionTok{:} \FloatTok{1.0}
\FunctionTok{\}}
\end{Highlighting}
\end{Shaded}
\end{itemize}

\subsubsection{6.2. Lancement et Analyse}\label{lancement-et-analyse}

\begin{Shaded}
\begin{Highlighting}[]
\ExtensionTok{mpirun} \AttributeTok{{-}np}\NormalTok{ 4 bin/covof2d cases/backward\_facing\_step.json}
\end{Highlighting}
\end{Shaded}

La longueur de recollement mesurée est \(x_r / H = 2.719\), en accord
direct avec la valeur expérimentale et numérique de Référence Académique
(\(x_r / H = 2.808\), écart \(< 3.2\%\)).

\begin{center}\rule{0.5\linewidth}{0.5pt}\end{center}

\subsection{7. Tutoriel 4 : Convection Naturelle Thermique de Vahl Davis
(Ra =
10⁴)}\label{tutoriel-4-convection-naturelle-thermique-de-vahl-davis-ra-10ux2074}

Ce cas étudie la convection thermique d'un fluide dans une cavité fermée
dont la paroi gauche est maintenue à \(T_h = 1.0\) et la paroi droite à
\(T_c = 0.0\), les parois supérieure et inférieure étant adiabatiques
(\(\partial T / \partial y = 0\)).

\subsubsection{7.1. Équations résolues}\label{uxe9quations-ruxe9solues}

\[\frac{\partial T}{\partial t} + \mathbf{u} \cdot \nabla T = \alpha \nabla^2 T\]
\[\mathbf{f}_{\text{boussinesq}} = \rho \mathbf{g} \beta (T - T_0)\]

\subsubsection{\texorpdfstring{7.2. Configuration
(\texttt{cases/natural\_convection\_ra1e4.json})}{7.2. Configuration (cases/natural\_convection\_ra1e4.json)}}\label{configuration-casesnatural_convection_ra1e4.json}

\begin{itemize}
\tightlist
\item
  \(Ra = \frac{\rho g \beta \Delta T L^3}{\mu \alpha} = 10^4\),
  \(Pr = 0.71\).
\item
  Calcul du nombre de Nusselt moyen à la paroi chaude :
  \[\overline{Nu} = \frac{1}{L_y} \int_0^{L_y} \left( -\frac{\partial T}{\partial x} \right)_{x=0} dy\]
\item
  \textbf{Résultat COVOF 2D} : \(\overline{Nu} = 2.7261\) (conforme à la
  plage benchmark \([2.2, 2.8]\)).
\end{itemize}

\begin{center}\rule{0.5\linewidth}{0.5pt}\end{center}

\subsection{8. Tutoriel 5 : Écoulement Diphasique VOF --- Rupture de
Barrage}\label{tutoriel-5-uxe9coulement-diphasique-vof-rupture-de-barrage}

L'effondrement gravitationnel d'une colonne d'eau (Martin \& Moyce) met
à l'épreuve le transport d'interface avec fortes déformations.

\begin{verbatim}
          t = 0.0 s                     t = 0.4 s (propagation front)
    +-----------------------+     +-----------------------+
    | ####                  |     |                       |
    | ####                  |     |                       |
    | ####                  |     | ##                    |
    | ####                  |     | #######======>        |
    +-----------------------+     +-----------------------+
\end{verbatim}

\subsubsection{\texorpdfstring{8.1. Configuration
(\texttt{cases/dam\_break.json})}{8.1. Configuration (cases/dam\_break.json)}}\label{configuration-casesdam_break.json}

\begin{itemize}
\tightlist
\item
  Colonne d'eau initiale : \(x \in [0, 0.2\text{ m}]\),
  \(y \in [0, 0.4\text{ m}]\), \(C = 1.0\).
\item
  Fluide 1 (Eau) : \(\rho_1 = 1000\text{ kg/m}^3\),
  \(\mu_1 = 0.001\text{ Pa}\cdot\text{s}\).
\item
  Fluide 2 (Air) : \(\rho_2 = 1.0\text{ kg/m}^3\),
  \(\mu_2 = 1.85 \times 10^{-5}\text{ Pa}\cdot\text{s}\).
\item
  Gravité : \(g_y = -9.81\text{ m/s}^2\).
\end{itemize}

\subsubsection{8.2. Lancement}\label{lancement-2}

\begin{Shaded}
\begin{Highlighting}[]
\ExtensionTok{mpirun} \AttributeTok{{-}np}\NormalTok{ 4 bin/covof2d cases/dam\_break.json}
\end{Highlighting}
\end{Shaded}

Le front d'onde atteint une vitesse maximale de \(2.01\text{ m/s}\),
caractéristique du régime d'impact visqueux.

\begin{center}\rule{0.5\linewidth}{0.5pt}\end{center}

\subsection{9. Tutoriel 6 : Tension de Surface et Saut de Pression de
Laplace}\label{tutoriel-6-tension-de-surface-et-saut-de-pression-de-laplace}

Ce cas valide le calcul de tension superficielle par la méthode CSF pour
une gouttelette d'eau statique de rayon \(R = 0.25\text{ m}\) au repos.

\subsubsection{9.1. Loi Théorique de
Laplace}\label{loi-thuxe9orique-de-laplace}

En 2D, le saut de pression à travers une interface circulaire d'énergie
de surface \(\sigma\) vaut :
\[\Delta P = P_{\text{int}} - P_{\text{ext}} = \frac{\sigma}{R}\] Pour
\(\sigma = 1.0\text{ N/m}\) et \(R = 0.25\text{ m}\) :
\[\Delta P_{\text{théorique}} = \frac{1.0}{0.25} = 4.000\text{ Pa}\]

\subsubsection{9.2. Lancement et
Validation}\label{lancement-et-validation}

\begin{Shaded}
\begin{Highlighting}[]
\ExtensionTok{mpirun} \AttributeTok{{-}np}\NormalTok{ 2 bin/covof2d cases/static\_droplet.json}
\end{Highlighting}
\end{Shaded}

La pression interne moyenne calculée par COVOF 2D donne
\(\Delta P = 3.86\text{ Pa}\) (erreur relative de \(3.38\%\)),
confirmant l'absence de courants parasites excessifs.

\begin{center}\rule{0.5\linewidth}{0.5pt}\end{center}

\subsection{10. Tutoriel 7 : Écoulement Autour d'un Cylindre par
Pénalisation de
Brinkman}\label{tutoriel-7-uxe9coulement-autour-dun-cylindre-par-puxe9nalisation-de-brinkman}

Ce tutoriel illustre le calcul d'un écoulement laminaire à \(Re=20\)
autour d'un cylindre de diamètre \(D=0.1\text{ m}\) placé à
\((x_c, y_c) = (0.5, 0.5)\).

\subsubsection{\texorpdfstring{10.1. Configuration
(\texttt{cases/obstacle\_cylinder\_re20.json})}{10.1. Configuration (cases/obstacle\_cylinder\_re20.json)}}\label{configuration-casesobstacle_cylinder_re20.json}

\begin{Shaded}
\begin{Highlighting}[]
\StringTok{"obstacle"}\ErrorTok{:} \FunctionTok{\{}
  \DataTypeTok{"enabled"}\FunctionTok{:} \KeywordTok{true}\FunctionTok{,}
  \DataTypeTok{"type"}\FunctionTok{:} \StringTok{"CYLINDER"}\FunctionTok{,}
  \DataTypeTok{"xc"}\FunctionTok{:} \FloatTok{0.5}\FunctionTok{,} \DataTypeTok{"yc"}\FunctionTok{:} \FloatTok{0.5}\FunctionTok{,}
  \DataTypeTok{"radius"}\FunctionTok{:} \FloatTok{0.05}
\FunctionTok{\}}
\end{Highlighting}
\end{Shaded}

\subsubsection{10.2. Lancement}\label{lancement-3}

\begin{Shaded}
\begin{Highlighting}[]
\ExtensionTok{mpirun} \AttributeTok{{-}np}\NormalTok{ 2 bin/covof2d cases/obstacle\_cylinder\_re20.json}
\end{Highlighting}
\end{Shaded}

Deux tourbillons symétriques stationnaires se développent en aval du
cylindre avec une longueur de recirculation \(L/D \approx 0.55-0.77\).

\begin{center}\rule{0.5\linewidth}{0.5pt}\end{center}

\subsection{11. Post-Traitement Avancé avec
ParaView}\label{post-traitement-avancuxe9-avec-paraview}

COVOF 2D exporte ses données au format rectilinéaire partitionné VTK : *
Fichier maître multi-processus : \texttt{output\_cas/cas\_XXXXXX.pvtr} *
Fichiers de sous-domaines locaux :
\texttt{output\_cas/cas\_p0\_XXXXXX.vtr}, \texttt{cas\_p1\_XXXXXX.vtr}

\subsubsection{11.1. Étapes sous
ParaView}\label{uxe9tapes-sous-paraview}

\begin{enumerate}
\def\labelenumi{\arabic{enumi}.}
\tightlist
\item
  Ouvrir le groupe de fichiers \texttt{cas\_..pvtr}.
\item
  Cliquer sur \textbf{Apply} dans le panneau \emph{Properties}.
\item
  \textbf{Afficher le champ de vitesse} : dans la barre déroulante de
  coloration, choisir \texttt{Velocity}.
\item
  \textbf{Tracer des lignes de courant (Streamlines)} :

  \begin{itemize}
  \tightlist
  \item
    Menu \emph{Filters} \(\to\) \emph{Common} \(\to\) \emph{Stream
    Tracer}.
  \item
    Dans \emph{Seeds}, choisir \emph{Line Source} ou \emph{Point Cloud}.
  \item
    Ajuster la résolution et cliquer sur \textbf{Apply}.
  \end{itemize}
\item
  \textbf{Visualiser l'interface diphasique} :

  \begin{itemize}
  \tightlist
  \item
    Appliquer le filtre \emph{Contour} sur le champ
    \texttt{ColorFunction} (fraction volumique \(C\)) à la valeur
    \texttt{0.5}.
  \end{itemize}
\item
  \textbf{Profils de vitesse 1D} :

  \begin{itemize}
  \tightlist
  \item
    Menu \emph{Filters} \(\to\) \emph{Data Analysis} \(\to\) \emph{Plot
    Over Line}.
  \item
    Spécifier les coordonnées de la ligne (ex: \((x, y) = (0.5, 0.0)\) à
    \((0.5, 1.0)\)) pour visualiser directement \(u(y)\).
  \end{itemize}
\end{enumerate}

\begin{center}\rule{0.5\linewidth}{0.5pt}\end{center}

\subsection{12. Automatisation Python et Études de
Convergence}\label{automatisation-python-et-uxe9tudes-de-convergence}

Pour analyser l'ordre de convergence d'un schéma spatial, COVOF 2D
fournit le script
\href{scripts/verify_cases.py}{\texttt{scripts/verify\_cases.py}}. Vous
pouvez écrire un script simple pour balayer plusieurs résolutions :

\begin{Shaded}
\begin{Highlighting}[]
\ImportTok{import}\NormalTok{ subprocess}
\ImportTok{import}\NormalTok{ json}
\ImportTok{import}\NormalTok{ numpy }\ImportTok{as}\NormalTok{ np}

\NormalTok{grids }\OperatorTok{=}\NormalTok{ [}\DecValTok{16}\NormalTok{, }\DecValTok{32}\NormalTok{, }\DecValTok{64}\NormalTok{, }\DecValTok{128}\NormalTok{]}
\NormalTok{errors }\OperatorTok{=}\NormalTok{ []}

\ControlFlowTok{for}\NormalTok{ N }\KeywordTok{in}\NormalTok{ grids:}
    \ControlFlowTok{with} \BuiltInTok{open}\NormalTok{(}\StringTok{"cases/poiseuille\_variable\_viscosity.json"}\NormalTok{, }\StringTok{"r"}\NormalTok{) }\ImportTok{as}\NormalTok{ f:}
\NormalTok{        cfg }\OperatorTok{=}\NormalTok{ json.load(f)}
\NormalTok{    cfg[}\StringTok{"mesh"}\NormalTok{][}\StringTok{"ny"}\NormalTok{] }\OperatorTok{=}\NormalTok{ N}
\NormalTok{    cfg[}\StringTok{"output"}\NormalTok{][}\StringTok{"output\_dir"}\NormalTok{] }\OperatorTok{=} \SpecialStringTok{f"output\_conv\_}\SpecialCharTok{\{}\NormalTok{N}\SpecialCharTok{\}}\SpecialStringTok{"}
    
    \ControlFlowTok{with} \BuiltInTok{open}\NormalTok{(}\StringTok{"cases/tmp\_case.json"}\NormalTok{, }\StringTok{"w"}\NormalTok{) }\ImportTok{as}\NormalTok{ f:}
\NormalTok{        json.dump(cfg, f)}
        
\NormalTok{    subprocess.run([}\StringTok{"mpirun"}\NormalTok{, }\StringTok{"{-}np"}\NormalTok{, }\StringTok{"4"}\NormalTok{, }\StringTok{"bin/covof2d"}\NormalTok{, }\StringTok{"cases/tmp\_case.json"}\NormalTok{])}
    \CommentTok{\# Calcul de l\textquotesingle{}erreur L1 par rapport à la solution exacte...}
\end{Highlighting}
\end{Shaded}

\begin{center}\rule{0.5\linewidth}{0.5pt}\end{center}

\subsection{13. Créer un Nouveau Cas Test à Partir de
Zéro}\label{cruxe9er-un-nouveau-cas-test-uxe0-partir-de-zuxe9ro}

Pour modéliser un nouveau problème CFD :

\begin{enumerate}
\def\labelenumi{\arabic{enumi}.}
\item
  \textbf{Créer un fichier JSON} dans
  \texttt{cases/mon\_nouveau\_cas.json}.
\item
  \textbf{Définir la géométrie} : dimensions physiques \(L_x, L_y\) et
  nombre de mailles \(n_x, n_y\).
\item
  \textbf{Sélectionner les propriétés fluides} : masse volumique
  \(\rho\) et viscosité dynamique \(\mu\).
\item
  \textbf{Poser les conditions aux limites} :

  \begin{itemize}
  \tightlist
  \item
    Paroi étanche immobile :
    \texttt{"type":\ "DIRICHLET",\ "u":\ 0.0,\ "v":\ 0.0}
  \item
    Entrée de canal :
    \texttt{"type":\ "DIRICHLET",\ "u":\ U\_0,\ "v":\ 0.0}
  \item
    Sortie libre convective : \texttt{"type":\ "NEUMANN"}
  \item
    Symétrie axiale : \texttt{"type":\ "FREESLIP"}
  \item
    Périodicité spatiale : \texttt{"type":\ "PERIODIC"}
  \end{itemize}
\item
  \textbf{Régler le solveur} : pour la plupart des cas laminaires,
  \texttt{r\_al:\ 100.0}, \texttt{uzawa\_iter:\ 1},
  \texttt{tol:\ 1.0e-7}, solveur \texttt{BICGSTAB2}.
\item
  \textbf{Lancer et visualiser} :

\begin{Shaded}
\begin{Highlighting}[]
\ExtensionTok{mpirun} \AttributeTok{{-}np}\NormalTok{ 4 bin/covof2d cases/mon\_nouveau\_cas.json}
\ExtensionTok{paraview}\NormalTok{ output\_mon\_nouveau\_cas/mon\_nouveau\_cas\_..pvtr}
\end{Highlighting}
\end{Shaded}
\end{enumerate}

\end{document}
